%------------------------------------------------------------------------------%
\section{Demonstração direta}

A demonstração direta é aquela que, por meio de uma sequência lógica, parte
das hipóteses e com argumentações coerentes concluímos a veracidade da tese.

%------------------------------------------------------------------------------%
\begin{proposicao}{Quadrado de número par}{prop:par}
  Se $x$ é par, então $x^2$ é par.
\end{proposicao}
%------------------------------------------------------------------------------%

%------------------------------------------------------------------------------%
\begin{proof}
  Essa proposição nos diz que se um número $x$ é par, então o seu quadrado
  também é par. Note que dizemos que um número inteiro é par, se ele tem o
  formato $2k$, com $k \in \Z$ e dizemos que um número inteiro é ímpar,
  se ele tem o formato $2k+1$, com $k \in \Z$. Vamos inicialmente
  destacar a hipótese e a tese dessa proposição:
  \[
    \text{Hipótese: } x \text{ é par}.
  \]
  \[
    \text{Tese: } x^2 \text{ é par}.
  \]
  Assim, partindo da tese, se $x$ é um número par, então $x=2k$,
  onde $k$ é algum número inteiro. Assim,
  \[
    x^2
    = xx
    = (2k)(2k)
    = 4k^2
    = 2(2k^2)
  \]
  Chamando $p=2k^2$, e notando que $2k^2$ é um número inteiro, temos que
  $x^2=2p$, então podemos concluir que $x^2$ é, também, um número par.
\end{proof}
%------------------------------------------------------------------------------%

Agora vamos provar o famoso Teorema de Pitágoras

%------------------------------------------------------------------------------%
\begin{teorema}{Teorema de Pitágora}{teo:teorema-pitagoras-novamente}
  Se $T$ é um triângulo retângulo
  de catetos $b$ e $c$ e hipotenusa $a$, então $a^2=b^2+c^2$.
\end{teorema}
%------------------------------------------------------------------------------%

%------------------------------------------------------------------------------%
\begin{figure}[h]
  \centering
  \includegraphics{teo-de-pitagoras}
  \caption{Triângulo $T$}
  \label{pita}
\end{figure}
%------------------------------------------------------------------------------%

%------------------------------------------------------------------------------%
\begin{figure}[h]
  \centering
  \includegraphics{teorema-de-pitagoras-quadrado}
  \caption{Quadrado formado pela junção de 4 triângulos retângulos iguais.}
  \label{pita1}
\end{figure}
%------------------------------------------------------------------------------%

%------------------------------------------------------------------------------%
\begin{proof}
  Aqui já sabemos que as hipótese são que $T$ é um triângulo retângulo de
  catetos $b$ e $c$ e hipotenusa $a$ e a tese, que é o resultado que queremos
  chegar é  $a^2=b^2+c^2$. Logo, considere um triângulo retângulo T conforme a
  Figura~\ref{pita}:

  A partir disto, vamos pegar 4 triângulos iguais os da Figura~\ref{pita} e
  posicionar estes triângulos de modo que formem um quadrado de lado $b+c$.
  Note que formamos também um quadrado interno de lados cujos lados são as
  hipotenusas $a$ dos triângulos, conforme ilustra a Figura~\ref{pita1}.

  Podemos calcular a área desse quadrado formado de duas formas. A primeira é
  notando que o lado desse quadrado é $b+c$, assim:
  \begin{equation}
    \label{eqpi}
    A_{\text{quadrado}} = (b+c)^2=b^2+2bc+c^2
  \end{equation}
  Por outro lado, podemos calcular a área desse quadrado maior somando a
  área dos 5 objetos que o compõem: 4 triângulos retângulos iguais e
  um quadrado menor de lado $a$.
  Ora,
  \begin{gather*}
    A_{\text{triangulo}} = \frac{bc}{2} \\
    A_{\text{quadrado menor}} = a^2
  \end{gather*}
  Assim,
  \begin{equation}\label{eqpi2}
    A_{\text{quadrado}}
    = 4\frac{bc}{2}+a^2
    = 2bc+a^2
  \end{equation}
  Igualando as expressões obtidas em (\ref{eqpi}) e (\ref{eqpi2})
  chegamos que
  \[
    b^2+2bc+c^2= 2bc+a^2
  \]
  organizando os termos e fazendo os cancelamentos necessários, obtemos
  \[
    a^2=b^2+c^2.
  \]
\end{proof}
%------------------------------------------------------------------------------%

